\documentclass{article}
\usepackage[utf8]{inputenc}
\usepackage{enumitem}
\usepackage{hyperref}
\title{Byggeveiledning for EMP-generator}
\author{}
\date{}

\begin{document}
\maketitle

\section{Sikkerhetsregler}
\begin{itemize}
\item Arbeid alltid på et ryddig, isolerende underlag.
\item Bruk vernebriller og isolerende hansker ved testing.
\item Ha en brannslukker tilgjengelig.
\item Utlad alltid kondensatoren før du berører kretsen.
\item Ikke rett spolen mot personer, dyr eller elektronikk du ikke har tenkt å ødelegge.
\end{itemize}

\section{Forberedelser}
Skaff alle komponenter iht. BOM. Verktøy: loddebolt, multimeter, oscilloskop (minst 100 MHz), høyspenningsprobe (100:1), vernebriller, isolerende hansker, varmepistol, superlim, 3D-printer (for kabinett og spoleform).

\section{PCB-montering}
\begin{enumerate}
\item Lodde komponentene på PCB i rekkefølge: motstander, dioder, MOSFET, LED, brytere, transformator.
\item Vær nøye med polaritet på dioder, LED og elektrolyttkondensator (C5, hvis brukt).
\item Transformator T1: lim ferrittkjernen til PCB, lodde viklingene direkte til pads.
\item Koble batteriklemmer og eksterne komponenter via skrueterminaler.
\end{enumerate}

\section{Transformatorvikling}
Se \texttt{transformer.tex} for detaljer.

\section{Gnistgap}
\begin{enumerate}
\item Kapp to 15 mm lange biter av 1.6 mm wolframelektrode. Slip endene flate med et slipestein eller sandpapir.
\item Monter dem i en holder av plexiglass med messingskruer slik at avstanden kan justeres. Bor hull for elektrodene og gjeng for skruene.
\item For trigger: bor et 1 mm hull midt mellom elektrodene. Lim inn en tynn, isolert ledning (0.2 mm emaljert) med 0.5 mm klaring til begge hovedelektroder. Denne ledningen er triggerelektroden.
\item Still inn gapet til ca. 0.5 mm ved hjelp av et blad (følerlære). Dette kan justeres senere under testing.
\end{enumerate}

\section{Spole}
\begin{enumerate}
\item 3D-print en sirkulær form med spiralrille (20 mm indre, 100 mm ytre, 2.5 mm dybde). Filen finnes i \texttt{enclosure/head.stl}.
\item Vikle 6 tørn 10 AWG tråd i rillen, lim med superlim underveis for å holde vindingene på plass.
\item La to ender stikke ut (minst 15 cm) for tilkobling. Fjern emaljen fra endene med sandpapir.
\end{enumerate}

\section{Sluttmontering}
\begin{enumerate}
\item Plasser PCB, batterier og gnistgap i håndtaket (3D-printet). Håndtaket bør ha utsparinger for brytere og LED.
\item Koble piezo-tenneren til trigger-elektroden. Den ene ledningen fra piezo går til triggerelektroden, den andre til jord (eller en av hovedelektrodene).
\item Koble hovedkondensator C1 mellom utgang (D1 katode) og jord. Bruk korte, tykke ledninger (minst 2.5 mm²) for å minimere induktans.
\item Koble gnistgapet mellom C1 pluss og spole. Spolens andre ende til jord.
\item Isoler alle høyspentforbindelser med krympestrømpe eller silikon.
\item Monter spolen foran på hodet, og fest hodet til håndtaket.
\end{enumerate}

\section{Første gangs oppstart}
\begin{enumerate}
\item Sjekk alle loddinger og tilkoblinger visuelt.
\item Uten batterier, mål motstand mellom høyspentutgang og jord – skal være uendelig (åpen krets).
\item Sett i batterier, slå på S3. Hold S1 inne og mål spenningen over C1 med et voltmeter (1000 V DC). Spenningen skal stige til 400–450 V i løpet av 5–10 sekunder. LED skal lyse.
\item Slipp S1. Spenningen skal holde seg. Trykk S2 for å lade ut. **Hold knappen inne i minst 5 sekunder.** Spenningen skal falle til under 10 V. Verifiser med voltmeter før du berører kretsen.
\item Test trigger: Lad opp til 400 V, trykk på piezo-knappen. Du skal høre et skarpt smell og se en gnist i gapet. Hvis ikke, juster gapet eller sjekk piezo-tilkobling.
\end{enumerate}
\end{document}
